\PassOptionsToPackage{table,xcdraw}{xcolor}
\documentclass[11pt, a4paper]{article}

\usepackage[T1]{fontenc}
\usepackage[utf8]{inputenc}
\usepackage{tgpagella}
\usepackage[spanish, es-noshorthands]{babel}
\usepackage{amsmath, amssymb}
\usepackage[most]{tcolorbox}
\usepackage{tabularx}
\usepackage{tikz}
\usepackage[american]{circuitikz}
\usetikzlibrary{shapes.geometric, arrows.meta, positioning, calc}
\usepackage[margin=2.5cm]{geometry}
\usepackage{fancyhdr}

\pagestyle{fancy}
\fancyhf{}
\setlength{\headheight}{16pt}
\lhead{\textbf{UCB Sede Tarija} | Ing. Mecatrónica}
\rhead{IMT-221 | Concepto High-Side PMOS (E4)}
\renewcommand{\headrulewidth}{0.4pt}
\definecolor{ucbblue}{RGB}{0, 51, 102}

\begin{document}

\begin{tcolorbox}[colback=orange!10, colframe=orange, title=\textbf{Análisis de Riesgo: El Componente Candidato IRF4905 (PMOS)[cite: 21]}]
    \textbf{Objetivo:} Demostrar por qué el componente aportado no puede usarse directamente con pines lógicos de $5\,\text{V}$ en una arquitectura High-Side[cite: 21]. Este circuito NO está autorizado para ensamble, es puro análisis lógico previo para E4[cite: 21].
\end{tcolorbox}

\section*{El Problema de Adaptación de Niveles Lógicos en PMOS[cite: 21]}

En una configuración High-Side con un transistor de canal P (PMOS), el Source se conecta típicamente al riel positivo de alimentación principal ($+12\,\text{V}$)[cite: 21]. El PMOS se enciende cuando su Gate es más negativo que su Source, es decir, cuando $\mathbf{V_{GS} < V_{GS(th)}}$ (un valor fuertemente negativo)[cite: 21].

\vspace{0.4cm}
\begin{center}
    \begin{circuitikz}[scale=1.1, transform shape, thick]
        % Source at +12V
        \draw (0,3) node[above]{$+12\,\text{V}$} -- (0,1.5) coordinate (source_node);
        
        % PMOS (pigfete) - El P-Channel enhancement tiene Source arriba por defecto
        \draw (0,0) node[pigfete] (M2) {};
        
        \node[right=0.1cm] at (M2.S) {S};
        \node[right=0.1cm] at (M2.D) {D};
        \node[above left=0.05cm] at (M2.G) {G};
        \node[below right=0.05cm] at (M2.B) {B};
        
        % Conexiones Source y Drain
        \draw (source_node) node[circle,fill,inner sep=1pt]{} -- (M2.S);
        \draw (M2.D) -- (0,-1) coordinate (drain_node) node[circle,fill,inner sep=1pt]{};
        
        % Conexión del Bulk a Source (por la derecha)
        \draw (M2.B) -- ++(0.6,0) |- (source_node);
        
        % Carga a tierra
        \draw (drain_node) to[R, l=Carga ($100\,\Omega$)] (0,-3) node[ground]{};
        
        % Control Gate (Señal Absoluta 0/5V)
        \draw (-3.5, 0) node[circle,fill,inner sep=1pt]{} node[above]{$V_G$} -- (M2.G);
        \draw (-3.5, 0) to[V, l_={\parbox{2.5cm}{\centering Señal Lógica\\$0/5\,\text{V}$ absolutos}}] (-3.5, -3) node[ground]{};
        
        % Marcador V_GS
        \draw[<->, red, thick] (-2.5, 0) -- (-2.5, 1.5) node[midway, left] {$V_{GS}$};
        \draw[dashed, red] (source_node) -- (-2.5, 1.5);
    \end{circuitikz}
\end{center}

\vspace{0.4cm}
\textbf{Análisis de Estados Lógicos (Falla Crítica de Apagado):}
\begin{enumerate}
    \item \textbf{Intento de OFF (Apagar el Transistor)[cite: 21]:} \\
    Usted envía un `1` lógico ($+5\,\text{V}$) desde su microcontrolador hacia el Gate ($V_G = 5\,\text{V}$ absolutos respecto a tierra)[cite: 21]. 
    Dado que el Source está amarrado al bus principal de $+12\,\text{V}$ ($V_S = 12\,\text{V}$):
    \begin{equation*}
        V_{GS} = V_G - V_S = 5\,\text{V} - 12\,\text{V} = \mathbf{-7\,\text{V}} \quad \text{[cite: 21]}
    \end{equation*}
    Con $V_{GS} = -7\,\text{V}$, el PMOS \textbf{seguirá fuertemente ENCENDIDO}. ¡Un GPIO directo de $0-5\,\text{V}$ no puede apagar un PMOS referenciado a $12\,\text{V}$![cite: 21] 
    
    \item \textbf{Conclusión de Diseño (Puerta E4):} Para apagarlo (Lograr $V_{GS} \approx 0\,\text{V}$), el controlador del Gate tendría que empujar el voltaje de compuerta hasta los $+12\,\text{V}$ completos[cite: 21]. Se hace mandatorio diseñar un circuito \textit{Level Shifter} y un Driver (Ej. BJT o Bootstrap) antes de poder operar un PMOS o NMOS en configuración High-Side de forma segura[cite: 21].
\end{enumerate}

\end{document}
